\chapter{Proving PRF-Induced MAC Security Equationally}
\lhead[\fancyplain{}{\bfseries\thepage}]{\fancyplain{}{\bfseries\rightmark}}

This chapter formalizes the security definition of pseudorandom-function–based message
authentication codes (PRF-induced MACs) introduced in Section~\ref{sec:computational-cryptography}. We consider the standard proof by reduction establishing the security of the MAC scheme $\mathrm{MAC}^F$, constructed from a pseudorandom function $F$, as outlined in that section.

The goal is to express this security definition and its proof within the $\lamBLL$
framework. To this end, we encode the components of the MAC-forgery experiment as
$\lamBLL$ terms and reconstruct the reduction argument using the equational reasoning
principles. In particular, we formalize the adversary, the MAC oracle, and the forgery experiment, and we prove the security of $\mathrm{MAC}^F$ by showing that any successful forgery induces a distinguisher for the underlying pseudorandom function.

First of all let's begin by recalling the MAC-forgery experiment $\macexp$ and its pseudocode, shown in Figure~\ref{fig:macforge-recall}.
\begin{figure}[htpb]
    \centering
    \begin{align*}
    &\macforge(n): \\
    &\quad k \gets Gen(1^n) \\
    &\quad (m,t) \gets A^{Mac_k(\cdot)}(1^n) \\
    &\quad \mathbb{Q} \gets \{m \ | \ A \ \text{queries} \ Mac_k(m) \} \\
    &\quad \mathtt{return}\,\,(m \notin \mathbb{Q} \land \vrfy_k(m,t) = 1)
    \end{align*}
    \caption{Pseudocode for $\mathtt{MacForge}$ experiment}
    \label{fig:macforge-recall}
\end{figure}

We now describe how the components of this experiment are typed in $\lamBLL$. The key-generation algorithm $Gen$ is a probabilistic computation that, given the
security parameter, produces a secret key of length polynomial in $i$. The adversary is modeled as a higher-order term that receives oracle access to the MAC
function and outputs a candidate forgery. Since the adversary may query the MAC oracle
a polynomial number $q(i)$ of times, the oracle is provided under a graded modality.

\begin{align*}
\vdash_v^\Theta &: \UU \multimap \SS[i] && \tag{Type for $Gen$} \\
\vdash_v^\Theta &: \vrfy : (\SS[i] \otimes \SS[i]) \multimap \BB && \tag{Type for $\vrfy$} \\
\vdash_v^\Theta &: Oracle : \bang_{q}\bigl(\SS[i] \multimap \SS[i]\bigr) && \tag{Type for $Oracle$} \\
\vdash_c^\Theta &: 
Adv : Oracle \multimap \bigl((\SS[i] \otimes \SS[i]) \otimes \SS[r]\bigr) && \tag{Type for $Adv$}
\end{align*}

Here, $i$, $r$, and $q$ are polynomials. The oracle type $(\SS[i] \multimap \SS[i])$ represents the MAC function mapping messages of length $i$ to tags of length $i$, wrapped in a graded modality to enforce a polynomial bound on queries. The adversary's output $(\SS[i] \otimes \SS[i]) \otimes \SS[r]$ contains a candidate message/tag pair and the ledger of all queries performed. The function $\vrfy$ checks tag validity, and $\mathtt{contains}$ ensures the adversary does not output a previously queried message. The type of the function $\mathtt{contains}$ is \[ \vdash_v^\Theta: \mathtt{contains} : \SS[r] \otimes \SS[i] \multimap \BB \]

The $\lamBLL$ term for the experiment is represented in Figure~\ref{fig:macforge-term}

\begin{figure}[ht]
    \centering
    \begin{align*}
    &MacForge \triangleq \\
    &\quad \letin{k = Gen *}{} \\
    &\quad \letin{o = Oracle\, k}{} \\
    &\quad \letin{\langle mt, s \rangle = Adv\, o}{} \\
    &\quad \letin{\langle m, t \rangle = mt}{} \\
    &\quad \letin{Q = s}{} \\
    &\quad \letin{b = \vrfy\, \langle m, t \rangle}{} \\
    &\quad \mathtt{if}\, b\, \mathtt{then}\, \mathtt{not}(\mathtt{contains}(Q, m)) \\
    &\quad \mathtt{else}\, \false
    \end{align*}
    \caption{$\lamBLL$ term for the $\mathtt{MacForge}$ experiment. The experiment returns $\true$ if the adversary outputs a fresh valid message/tag pair and $\false$ otherwise.}
    \label{fig:macforge-term}
\end{figure}

In this formulation the algorithm $Gen$ produces a random key of length $i$. The $Oracle\, k$ is the MAC function under the key $k$, duplicated according to the polynomial $q(i)$. $Adv$ is a computation interacting with the oracle, producing a candidate message/tag pair and the query ledger $Q$. The algorithm $\vrfy$ checks the correctness of the MAC tag for the candidate message. $\mathtt{contains}$ checks whether the candidate message is already present in the ledger. The experiment returns $\true$ if and only if the tag is valid and the message was not queried previously, otherwise it returns $\false$.